\documentclass[12pt]{article}
\usepackage[T1]{fontenc}
\usepackage[utf8]{inputenc}
\usepackage{lmodern}
\usepackage{textcomp}
\usepackage{enumitem}
\usepackage{geometry}
\geometry{margin=1in}
\begin{document}
\begin{center}
Answer Key - Political Science - Undergraduate Level Assessment 5 \
\small (Answers are ordered exactly as in the question paper.)
\end{center}

\section*{Multiple Choice}
\begin{enumerate}[label=\arabic*.]
\item \textbf{Answer:} A
\\ \textit{Options:} A) Married white male; B) Unmarried white female; C) African American, male or female; D) Youths under the age of 25, male or female
\\ \textit{Note:} Married white males are statistically more likely to align with Republican values due to cultural, economic, and social factors, including traditional views on family and a preference for conservative fiscal policies.
\item \textbf{Answer:} B
\\ \textit{Options:} A) Political policies designed to alter demographic characteristics of a state.; B) All of these options.; C) Legislation that criminalises certain cultural behaviours or practice.; D) The use of military force to conduct ethnic-cleansing through displacement and killing.
\\ \textit{Note:} The correct answer is "All of these options" because societal security encompasses a wide range of threats, including economic instability, political violence, and social unrest, all of which can undermine the stability and cohesion of a society.
\item \textbf{Answer:} C
\\ \textit{Options:} A) A news story focuses on a politician's scandals rather than achievements.; B) A television news anchor reports an event before the station's rivals.; C) A reporter announces which candidate leads in a public opinion poll.; D) A newspaper editor prints stories about long-term political developments.
\\ \textit{Note:} The correct answer exemplifies "horse-race journalism" because it focuses on the competition between candidates by emphasizing who is leading in polls rather than discussing their policies or qualifications, creating a narrative centered on the race itself.
\item \textbf{Answer:} B
\\ \textit{Options:} A) The death penalty may only be imposed upon citizens.; B) Under some circumstances, the death penalty may not violate the Eighth Amendment.; C) Under some circumstances, the death penalty may violate the Third Amendment.; D) States may execute any adult regardless of intellectual disability.
\\ \textit{Note:} The Supreme Court has ruled that the death penalty can be constitutional under certain conditions, provided it does not constitute cruel and unusual punishment as prohibited by the Eighth Amendment.
\item \textbf{Answer:} D
\\ \textit{Options:} A) Benevolent hegemony; B) Ineffectiveness of multilateralism; C) American power; D) All of the above
\\ \textit{Note:} The correct answer "All of the above" is justified because each point listed in the options provides distinct yet complementary evidence that underscores the rationale and strategic advantages of American unilateralism in foreign policy.
\end{enumerate}

\section*{True / False}
\begin{enumerate}[label=\arabic*.]
\setcounter{enumi}{5}
\item \textbf{Answer:} True
\\ \textit{Options:} A) True; B) False
\\ \textit{Note:} A motion for cloture is a procedure in the Senate that allows senators to end a filibuster and proceed to a vote on a bill, requiring a supermajority of 60 votes to invoke.
\item \textbf{Answer:} True
\\ \textit{Options:} A) True; B) False
\\ \textit{Note:} The principle of checks and balances is enshrined in the Constitution to create a system where each branch of government-executive, legislative, and judicial-has the ability to limit the powers of the others, thereby preventing any one branch from gaining excessive power.
\item \textbf{Answer:} False
\\ \textit{Options:} A) True; B) False
\\ \textit{Note:} The statement is false because the U.S. Constitution grants Congress the authority to declare war, requiring the president to seek congressional approval for such actions.
\item \textbf{Answer:} False
\\ \textit{Options:} A) True; B) False
\\ \textit{Note:} The statement is false because in an open primary, registered independents are allowed to vote in the primary elections of any party they choose, regardless of their independent status.
\item \textbf{Answer:} False
\\ \textit{Options:} A) True; B) False
\\ \textit{Note:} The statement is false because oil has historically been a significant factor in shaping U.S. foreign policy, influencing military interventions, alliances, and economic strategies due to its critical role in global energy security and the economy.
\end{enumerate}

\section*{Long Form Response}
\begin{enumerate}[label=\arabic*.]
\setcounter{enumi}{10}
\item \textbf{Answer:} Revisionist historians argue that domestic concerns in the United States played a crucial role in the onset of the Cold War, suggesting that internal political pressures, economic interests, and cultural anxieties drove U.S. foreign policy decisions. The fear of communism infiltrating American society fueled a defensive posture against the Soviet Union, leading to the pursuit of aggressive containment strategies. This domestic paranoia was exacerbated by events such as the Red Scare, which heightened anti-communist sentiment and justified interventionist policies abroad. Consequently, these domestic issues not only shaped U.S. actions during the Cold War but also influenced international relations by creating a polarized world where ideological battles overshadowed diplomatic discourse, leading to prolonged tensions and conflicts. Thus, the interplay between domestic pressures and foreign policy decisions fundamentally shaped the trajectory of the Cold War era.
\\ \textit{Note:} correct because it effectively highlights how domestic political pressures, economic interests, and cultural anxieties in the United States shaped foreign policy decisions during the Cold War, illustrating the interplay between internal concerns and international relations.
\item \textbf{Answer:} Post-structuralist critiques play a crucial role in security studies by emphasizing the fluidity of power dynamics, the importance of identity, and the construction of security narratives. They challenge traditional views that regard security as a fixed state, instead arguing that security is socially constructed, shaped by discourses and power relations. This perspective allows for a deeper analysis of how identities-such as national, ethnic, or gender identities-are intertwined with security issues and how these identities influence perceptions of threat. Furthermore, post-structuralism highlights the narratives surrounding security, revealing how they can be manipulated to justify policies and actions. Overall, these critiques enrich our understanding of security by illuminating the complexities of power and identity in a global context.
\\ \textit{Note:} correct because it accurately identifies post-structuralist critiques as key to understanding the fluidity of power dynamics, the role of identity in shaping security, and the idea that security is a socially constructed narrative influenced by discourse and power relations.
\end{enumerate}

\vfill
\noindent\textit{End of Answer Key}
\end{document}